\section{Instant Contact Dynamics Experiment}

\subsection{Purpose}

The instant contact dynamics experiment attempted to remove
\texttt{first\_contact\_velocity} from the active dynamics path.  The motivation
was to make the collision rule more instantaneous, more geometric, and easier to
port to GLSL.  The intended model was:

\begin{itemize}
    \item use the current source and target velocities at the current step;
    \item compute the current relative normal velocity;
    \item store incoming normal momentum while the contact is compressing;
    \item return stored normal momentum while the contact is releasing;
    \item keep prior velocity history as diagnostic/reporting data only.
\end{itemize}

\subsection{Code Changes Made}

The experiment added a new runtime flag:

\begin{verbatim}
instant_contact_dynamics
\end{verbatim}

The flag was added to \texttt{Base} and read from \texttt{RUN\_CONFIGURATION} in
\texttt{gPCD}.  Many configuration files were given an
\texttt{instant\_contact\_dynamics} entry so the experimental path could be
enabled or disabled per configuration.

The normal source-owned path:

\begin{verbatim}
resolve_source_contact_momentum(...)
\end{verbatim}

was changed so that, when the flag was true, it routed to:

\begin{verbatim}
resolve_instant_source_contact_momentum(...)
\end{verbatim}

Wall contact prediction was also changed so that wall ghosts used the same
instant resolver when the flag was enabled.

\subsection{Instant Resolver Algorithm}

The new resolver used the current source and target velocities:

\begin{verbatim}
rel_v = target_velocity - source_velocity
rel_vn = dot(rel_v, contact_normal)
\end{verbatim}

If the contact was accumulating and the current relative normal velocity was not
closing, the resolver returned no velocity change.

The resolver then computed the zero-contact momentum using the existing
geometric zero-velocity contact state:

\begin{verbatim}
p_zero = geo_zero_contact_momentum(...)
\end{verbatim}

For multi-contact cases, it kept the existing source-owned area weighting:

\begin{verbatim}
weighted p_zero = available source momentum * contact_area / total_source_area
\end{verbatim}

The area-scaled contact momentum was computed from current overlap:

\begin{verbatim}
p_area = geo_area_scaled_contact_momentum(overlap_area, p_zero, contact_state)
\end{verbatim}

The experiment introduced an instant stored momentum field:

\begin{verbatim}
instant_stored_internal_normal_momentum
\end{verbatim}

During compression, the target stored value was:

\begin{verbatim}
target_stored = max(stored, p_area)
store_delta = target_stored - stored
\end{verbatim}

The actual transfer was capped by current closing momentum:

\begin{verbatim}
closing_momentum = reduced_mass * max(0, -rel_vn)
transfer = min(store_delta, closing_momentum)
stored += transfer
\end{verbatim}

During return, the resolver attempted to reduce stored momentum and convert it
back into source velocity.

\subsection{Additional Release Attempts}

The first wall tests showed that wall rebound could become infinitely slow.  The
contact kept accumulating because the current closing momentum became smaller
each frame, approaching zero without reliably entering release.

To address this, the experiment added a rule: if compression consumed all
available closing momentum, mark the contact as an instantaneous zero-velocity
point:

\begin{verbatim}
instant_zero_velocity_reached = true
internal_momentum_phase = returning
neo_motion_mode = releasing
geo_zero_velocity_source = instant_velocity_zero
\end{verbatim}

The experiment also added a pending-release field:

\begin{verbatim}
instant_release_pending_normal_momentum
\end{verbatim}

The goal was to release stored momentum even when the particle stopped exactly
at the zero point and the overlap area did not immediately shrink.

\subsection{Observed Problems}

\subsubsection{Wall Rebound Became Too Slow}

In \texttt{TwoParticleHorizontal}, particle-wall contact stored momentum but the
normal velocity approached zero slowly.  The contact remained in compression or
failed to release enough stored momentum.  This showed that the simple
instantaneous compression cap:

\begin{verbatim}
transfer = min(store_delta, closing_momentum)
\end{verbatim}

can drain velocity without giving the contact a robust return driver.

\subsubsection{Stopped Contact Did Not Release}

In \texttt{OneParticleBoundaryTest}, the wall contact reached about
\texttt{0.05} stored internal momentum, which matched the expected incoming
momentum.  However, once the particle stopped, release did not occur.  The
contact reported \texttt{returning/releasing}, but \texttt{irel} stayed zero and
velocity stayed zero.

This revealed that the release rule was still too dependent on geometry changing
after the zero point.  If the particle is exactly stopped, geometry does not
advance by itself, so release must have a clear driver.

\subsubsection{Release State Was Set Before Resolver Entry}

Another complication was that \texttt{update\_wall\_contact\_state()} could mark
the contact as \texttt{returning} before the resolver ran.  A transition test
inside the resolver was therefore unreliable: by the time the resolver saw the
state, the transition had already happened.

\subsection{Why The Experiment Was Rolled Back}

The instant dynamics path did not perform as well as the previous checked-in
model.  Particle-particle rebound and wall collisions behaved better without the
experimental branch.  The main unresolved issue is that removing velocity
history requires a replacement return rule that is strong enough to return
stored momentum, but not so strong that it becomes another hidden heuristic.

\subsection{Open Questions}

\begin{itemize}
    \item What is the correct purely geometric return driver once a contact is
    stopped at the zero point?
    \item Should stored internal normal momentum be returned according to change
    in overlap area, current relative velocity, or a separate contact state?
    \item Can the release rule be expressed without first-contact velocity while
    still preserving momentum in chain collisions?
    \item How should a fixed wall ghost differ from a moving real particle, if
    both must pass through the same source-owned resolver?
\end{itemize}

\subsection{Rollback Note}

After documenting this experiment, the Python and configuration changes from the
instant contact dynamics branch were restored to the previous checked-in state.
